\documentclass[11pt]{article}
\usepackage[a4paper, margin=2cm]{geometry}
\usepackage[utf8]{inputenc}
\usepackage[T1]{fontenc}
\usepackage[polish]{babel}
\usepackage{amsmath, amssymb, amsthm}
\usepackage{booktabs}
\usepackage{tabularx}
\usepackage{xcolor}
\usepackage{hyperref}
\usepackage{lmodern}
\usepackage[scale=0.85]{FiraMono}
\usepackage[scaled=0.95]{helvet}

\definecolor{docblue}{RGB}{0, 102, 204}

\hypersetup{
    colorlinks=true,
    linkcolor=docblue,
    citecolor=docblue,
    urlcolor=docblue
}

\title{\textbf{Dokumentacja Projektowa}\\Wzorowy System Biletowy}
\author{Jan Pulkowski, Miłosz Zajdel}
\date{\today}

\begin{document}

\maketitle
\tableofcontents
\newpage

\section{Wprowadzenie}
Wzorowy System Biletowy to aplikacja serwerowa realizująca proces rezerwacji i sprzedaży biletów.
System został zaimplementowany w języku Java 21 z wykorzystaniem środowiska Spring Boot oraz bazy danych SQLite.

\subsection{Quickstart}
Wszystkie zadania budowania realizuje skrypt \texttt{gradlew} (lub w systemie Windows \texttt{gradlew.bat}).
\begin{itemize}
  \item Uruchomienie -- \texttt{./gradlew~run}
  \item Kompilacja -- \texttt{./gradlew~build}
  \item Kompilacja dokumentacji z tex do pdf -- \texttt{./gradlew~buildDocs}
\end{itemize}

\subsection{Architektura}
Wzorowy system biletowy to monolityczny REST service.
\begin{itemize}
    \item \textbf{Baza Danych}: SQLite wykorzystywana w trybie wbudowanym, klucze eksponowane przez API, zatem w formacie UUID. W celach ułatwienia testów i pracy baza jest tworzona na nowo przy każdym uruchomieniu aplikacji.
    Zachowanie to można zmienić w pliku konfiguracyjnym \texttt{resources/application.yaml} -- klucz \texttt{spring.sql.init.mode}.
    \item \textbf{API}: Osobne endpointy dla każdego z modułów -- katalog ofert, procesor rezerwacji, procesor płatności. Implementacja według wzoru Kontroler-Usługa-Repozytorium.
\end{itemize}

\section{Katalog ofert}

\subsection{Model Oferty}

\begin{figure}[H]
\label{schema:offer}
Jeżeli typ pola nie jest podany to jest to string.
\begin{verbatim}
{
  "id": "UUID oferty",
  "name": "Nazwa oferty",
  "vendor": {
    "name": "Nazwa sprzedawcy"
  },
  "price": {
    "amount": "Cena oferty w formacie tekstowym",
    "currency": "Kod waluty"
  },
  "description": "Opis oferty",
  "quantity": "Liczba możliwych sprzedaży oferty -- liczba całkowita"
}
\end{verbatim}
\caption{Obiekt Offer}
\end{figure}


\subsection{Reprezentacja w bazie danych}

\begin{table}[ht]
  \centering
  \renewcommand{\arraystretch}{1.3} % Adds vertical padding to rows for readability
  \begin{tabularx}{\textwidth}{|l|l|l|X|l|}
    \hline
    \textbf{Kolumna} & \textbf{Typ SQL} & \textbf{Ograniczenia} & \textbf{Opis} & \textbf{Odpowiednik API} \\
    \hline
    id & TEXT & PRIMARY KEY & UUID oferty & \texttt{id} \\
    \hline
    name & TEXT & NOT NULL & Nazwa oferty & \texttt{name} \\
    \hline
    vendor\_id & INTEGER & NOT NULL, FK & Klucz zewnętrzny do tabeli \texttt{vendor} & Niedostępne przez api \\
    \hline
    currency & TEXT & NOT NULL & Kod waluty & \texttt{price.currency} \\
    \hline
    amount & TEXT & NOT NULL & Cena oferty w formacie tekstowym & \texttt{price.amount} \\
    \hline
    description & TEXT & NOT NULL & Opis oferty & \texttt{description} \\
    \hline
    quantity & INTEGER & NOT NULL, $\ge 0$ & Dostępność oferty & \texttt{quantity} \\
    \hline
  \end{tabularx}
  \caption{Struktura tabeli \texttt{offer}}
  \label{tab:schema-offer}
\end{table}

\begin{table}[ht]
  \centering
  \renewcommand{\arraystretch}{1.3}
  \begin{tabularx}{\textwidth}{|l|l|l|X|l|}
    \hline
    \textbf{Kolumna} & \textbf{Typ SQL} & \textbf{Ograniczenia} & \textbf{Opis} & \textbf{Odpowiednik API} \\
    \hline
    id & INTEGER & PK, AUTOINCREMENT & Klucz pierwotny & Niedostępne przez api \\
    \hline
    name & TEXT & NOT NULL & Unikalna indeksowana nazwa sprzedawcy & \texttt{vendor.name} \\
    \hline
  \end{tabularx}
  \caption{Struktura tabeli \texttt{vendor}}
  \label{tab:schema-vendor}
\end{table}

\newpage

\subsection{Endpointy API}

Ścieżka bazowa API katalogu ofert to \texttt{/offer}. Jeżeli pole nie jest podane jako opcjonalne to jest obowiązkowe i musi zostać wypełnione. Wszystkie zapytania zwracają \hyperref[schema:offer]{obiekt Offer}, jeżeli zapytanie się powiodło, w przeciwnym wypadku kod błędu.


\subsubsection{Endpoint \texttt{/\{id\}}}

\begin{itemize}
  \item Metoda: \texttt{GET}. Brak parametrów i ciała zapytania. Zwraca ofertę o UUID równym podanemu \texttt{id}.
  \item Metoda: \texttt{DELETE}. Brak parametrów i ciała zapytania. Usuwa i zwraca ofertę o UUID równym podanemu \texttt{id}.
\end{itemize}

\subsubsection{Endpoint \texttt{/search}}

Metoda: \texttt{GET}. Parametry: \texttt{vendor}, \texttt{name}. Brak ciała zapytania. Zwraca oferty pasujące nazwą i/lub nazwą sprzedawcy do parametrów zapytania. Obowiązkowe jest podanie co najmniej 1 parametru.

\subsubsection{Endpoint \texttt{/create}}

Metoda: \texttt{POST}. Brak parametrów. Ciało zapytania: \hyperref[schema:offer]{obiekt Offer}. Tworzy zadaną ofertę. Pole \texttt{id} jest nieobowiązkowe i ignorowane; pole \texttt{description} może być pustym stringiem. Zwraca utworzoną ofertę.

\subsubsection{Endpoint \texttt{/update}}

Metoda: \texttt{POST}. Brak parametrów. Ciało zapytania: \hyperref[schema:offer]{obiekt Offer}. Aktualizuje zadaną ofertę. Wszystkie pola poza \texttt{id} są nieobowiązkowe; pole \texttt{vendor} nie musi zostać wypełnione, jeżeli podano pole \texttt{vendor} to musi się ono zgadzać ze stanem dotychczasowym -- nie można go zmieniać. Oferta zostanie zaktualizowana o wszystkie pozostałe pola wypełnione w zapytaniu. Zwraca zaktualizowaną ofertę.

\section{Procesor rezerwacji}

\subsection{Model Rezerwacji}

\begin{figure}[H]
Jeżeli typ pola nie jest podany to jest to string.
\begin{verbatim}
{
  "id": "UUID klienta",
  "name": "Imię i nazwisko lub nazwa firmy",
  "email": "Adres e-mail klienta"
}
\end{verbatim}
\caption{Obiekt Client}
\label{schema:client}
\end{figure}

\begin{figure}[H]
Jeżeli typ pola nie jest podany to jest to string.
\begin{verbatim}
{
  "id": "UUID rezerwacji",
  "items": [
    "Zrzut stanu i cen pozycji koszyka -- lista obiektów OrderItem"
  ],
  "createTimestamp": "Znacznik czasu utworzenia",
  "lastModifiedTimestamp": "Znacznik ostatniej modyfikacji"
}
\end{verbatim}
\caption{Obiekt Reservation}
\label{schema:reservation}
\end{figure}

\begin{figure}[H]
Jeżeli typ pola nie jest podany to jest to string.
\begin{verbatim}
{
  "lineId": "UUID pozycji koszyka",
  "offerSnapshot": "zrzut oferty w momencie rezerwacji -- obiekt Offer",
  "quantity": "Liczba zarezerwowanych sztuk w danej pozycji -- liczba całkowita"
}
\end{verbatim}
\caption{Obiekt OrderItem}
\label{schema:order-item}
\end{figure}


\subsection{Reprezentacja w bazie danych}
\newpage
\begin{table}[t]
  \centering
  \renewcommand{\arraystretch}{1.3}
  \begin{tabularx}{\textwidth}{|l|l|l|X|l|}
    \hline
    \textbf{Kolumna} & \textbf{Typ SQL} & \textbf{Ograniczenia} & \textbf{Opis} & \textbf{Odpowiednik API} \\
    \hline
    id & TEXT & PRIMARY KEY & UUID klienta & \texttt{id} \\
    \hline
    name & TEXT & NOT NULL & Imię i nazwisko lub nazwa firmy & \texttt{name} \\
    \hline
    email & TEXT & NOT NULL & Adres e-mail klienta & \texttt{email} \\
    \hline
  \end{tabularx}
  \caption{Struktura tabeli \texttt{client}}
  \label{tab:schema-client}
\end{table}

\begin{table}[ht]
  \centering
  \renewcommand{\arraystretch}{1.3}
  \begin{tabularx}{\textwidth}{|l|l|l|X|l|}
    \hline
    \textbf{Kolumna} & \textbf{Typ SQL} & \textbf{Ograniczenia} & \textbf{Opis} & \textbf{Odpowiednik API} \\
    \hline
    id & TEXT & PRIMARY KEY & UUID rezerwacji & \texttt{id} \\
    \hline
    version & INTEGER & NOT NULL & Licznik modyfikacji & Niedostępne przez api \\
    \hline
    client\_id & TEXT & NOT NULL, FK & Klucz zewnętrzny do tabeli \texttt{client} & Niedostępne przez api \\
    \hline
    offer\_items\_data & TEXT & NOT NULL & Zrzut JSON \hyperref[schema:order-item]{obiektów OrderItem} & \texttt{items} \\
    \hline
    create\_timestamp & INTEGER & NOT NULL & Znacznik czasu utworzenia & \texttt{createTimestamp} \\
    \hline
    last\_modified\_timestamp & INTEGER & NOT NULL & Znacznik ostatniej modyfikacji & \texttt{lastModifiedTimestamp} \\
    \hline
  \end{tabularx}
  \caption{Struktura tabeli \texttt{reservation}}
  \label{tab:schema-reservation}
\end{table}

\newpage


\section{Procesor płatności}

\subsection{Model płatności}

\begin{figure}[H]
Jeżeli typ pola nie jest podany to jest to string.
\begin{verbatim}
{
  "id": "UUID płatności",
  "reservationId": "UUID opłacanej rezerwacji",
  "amount": "Kwota transakcji w formacie tekstowym",
  "currency": "Kod waluty",
  "paymentTimestamp": "Czas rejestracji wpłaty"
}
\end{verbatim}
\caption{Obiekt Payment}
\label{schema:payment}
\end{figure}


\subsection{Reprezentacja w bazie danych}

\begin{table}[hbt]
  \centering
  \renewcommand{\arraystretch}{1.3}
  \begin{tabularx}{\textwidth}{|l|l|l|X|l|}
    \hline
    \textbf{Kolumna} & \textbf{Typ SQL} & \textbf{Ograniczenia} & \textbf{Opis} & \textbf{Odpowiednik API} \\
    \hline
    id & TEXT & PRIMARY KEY & UUID płatności & \texttt{id} \\
    \hline
    reservation\_id & TEXT & NOT NULL, FK & Klucz zewnętrzny do tabeli \texttt{reservation} & \texttt{reservationId} \\
    \hline
    amount & TEXT & NOT NULL & Kwota transakcji w formacie tekstowym & \texttt{amount} \\
    \hline
    currency & TEXT & NOT NULL & Trzyliterowy kod waluty & \texttt{currency} \\
    \hline
    payment\_timestamp & INTEGER & NOT NULL & Czas rejestracji wpłaty & \texttt{paymentTimestamp} \\
    \hline
  \end{tabularx}
  \caption{Struktura tabeli \texttt{payment}}
  \label{tab:schema-payment}
\end{table}



\end{document}
